\section*{Capítulo \ref{ch:b2:isotope-geochemistry} --- Geoquímica isotópica}

\begin{solution}{exo:b2:isotope-geochemistry:1}
Tristán: $(0.51249/0.512638 - 1)\times10^4 = -2.9$. Mangaia: $(0.51285/0.512638 -
1)\times10^4 = +4.1$. MORB: $+8.5$. Las fuentes de Mangaia y del MORB, con
$\varepsilon_{\mathrm{Nd}}$ positiva, han tenido de media una Sm/Nd más alta que las \omterm{def:b2:planet-formation-meteorites:chondrule}{condritas}; la
fuente de Tristán, una más baja.
\end{solution}

\begin{solution}{exo:b2:isotope-geochemistry:2}
$\delta^{18}\mathrm{O}_{\mathrm{VPDB}} = (\num{2.0672e-3}/\num{2.0052e-3} - 1)\times1000 =
+30.9$ por mil en la escala \omterm{def:b2:isotope-geochemistry:standards}{VSMOW}. La calcita tiene $^{18}\mathrm{O}/^{16}\mathrm{O} =
\num{2.0672e-3}\times(1 - 0.0020) = \num{2.0631e-3}$, de donde $\delta^{18}\mathrm{O} =
(\num{2.0631e-3}/\num{2.0052e-3} - 1)\times1000 = +28.9$ por mil (\omterm{def:b2:isotope-geochemistry:standards}{VSMOW}).
\end{solution}

\begin{solution}{exo:b2:isotope-geochemistry:3}
Hace \qty{1.0}{Ga}: $e^{\lambda t} - 1 = 0.006545$, y $0.512638 - 0.1966\times0.006545 =
0.511351$. Hace \qty{4.567}{Ga}: $e^{\lambda t}-1 = 0.030243$, y $0.512638 -
0.1966\times0.030243 = 0.506692$.
\end{solution}

\begin{solution}{exo:b2:isotope-geochemistry:4}
$\delta\mathrm{D} = 8\times(-18.0) + 10 = -134$ por mil. Para el lago, $d = -42 -
8\times(-4.0) = -10$ por mil, por debajo del 10 de la recta: el agua queda por debajo de la \omterm{prop:b2:isotope-geochemistry:mwl}{recta de agua
meteórica}, como les ocurre a las aguas superficiales evaporadas (los isótopos ligeros se fueron primero).
\end{solution}

\begin{solution}{exo:b2:isotope-geochemistry:5}
$Q = 10^4\times\num{6.524e-12}\times0.1966/0.512638 = \num{2.502e-8}$ por año, o
\qty{25.0}{\per\giga\year}. El granito tiene $\mathrm{Sm}/\mathrm{Nd} = 6.0/30 = 0.200$,
frente a $0.145/0.457 = 0.317$ en CI: $f = 0.200/0.317 - 1 = -0.37$. En
\qty{0.5}{Gyr}: $\Delta\varepsilon = 25.0\times(-0.37)\times0.5 = -4.6$ unidades.
\end{solution}

\begin{solution}{exo:b2:isotope-geochemistry:6}
$\varepsilon \approx 3.5 + 25.0\times(-0.45)\times2.7 = 3.5 - 30.4 = -26.9$, frente a
$-27.3$: la aproximación es 0.4 unidades demasiado alta (no lo bastante negativa). La mayor parte de la
diferencia, unas 0.3 unidades, procede de $e^{\lambda t} > 1$: el padre, y por tanto el
déficit de crecimiento respecto al CHUR, era más abundante en el pasado que hoy
(\cref{exo:b2:isotope-geochemistry:10}); el resto procede del redondeo de $f$ y del valor inicial.
\end{solution}

\begin{solution}{exo:b2:isotope-geochemistry:7}
Con una fracción másica $x$ de granito, $\mathrm{Sr} = 320x + 130(1-x)$ y
$R = (320x\times0.7200 + 130(1-x)\times0.7030)/\mathrm{Sr}$. Para $x = 0.05$: Sr
\qty{139.5}{ppm}, $R = 0.70495$; $x = 0.10$: 149.0, 0.70665; $x = 0.20$: 168.0,
0.70948. Por la demostración de la proposición de mezcla, $R\,\mathrm{Sr}$ es lineal en Sr:
$R = 0.731632 - 3.7221/\mathrm{Sr}$, una recta de pendiente $-3.722$ (ppm) y
ordenada en el origen 0.7316 frente a $1/\mathrm{Sr}$. Comprobación en $x = 0.05$: $0.731632 - 3.7221/139.5
= 0.70495$.
\end{solution}

\begin{solution}{exo:b2:isotope-geochemistry:8}
A \qty{20}{\degreeCelsius}: $\delta^{18}\mathrm{O} = 1000/1.0098 - 1000 = -9.70$ por
mil, $\delta\mathrm{D} = 1000/1.084 - 1000 = -77.5$ por mil, pendiente $77.5/9.70 = 7.99$.
A \qty{0}{\degreeCelsius}: $-11.56$ y $-99.9$ por mil, pendiente 8.64. Ambas están cerca de
la pendiente 8 de la \omterm{prop:b2:isotope-geochemistry:mwl}{recta de agua meteórica}: el fraccionamiento de equilibrio del agua
entre líquido y vapor basta por sí solo para fijar cerca de 8 la razón entre los desplazamientos del deuterio y del
oxígeno-18.
\end{solution}

\begin{solution}{exo:b2:isotope-geochemistry:9}
La corteza contiene $0.0051\times27/1.25 = 0.110$, el 11 por ciento del neodimio. Balance de
masas: $0.110\times(-15) + 0.890\,\varepsilon_m = 0$, de modo que $\varepsilon_m = 1.65/0.890 =
+1.9$. Con un reservorio empobrecido que contiene una proporción $x$ a $+8.5$ y el resto a 0:
$0.110\times(-15) + 8.5x = 0$, $x = 1.65/8.5 = 0.19$: cerca de una quinta parte del
neodimio.
\end{solution}

\begin{solution}{exo:b2:isotope-geochemistry:10}
El ritmo de crecimiento es $\lambda p\,e^{\lambda t}$, con $e^{\lambda t} \approx 1 + \lambda t$;
promediado sobre el intervalo que va de la separación $T$ al presente, el factor es
$1 + \lambda T/2$, de modo que los términos despreciados tienen un tamaño relativo $\lambda T/2$. A lo largo de la
edad de la Tierra, $\num{6.524e-12}\times\num{4.567e9}/2 = 0.015$, un 1.5 por ciento. Para la
corteza del \cref{exo:b2:isotope-geochemistry:6}, $\lambda T/2 = 0.0088$, y el error es $0.0088\times30.4 = 0.27$
unidades. El lento cambio de la razón del CHUR por la que se divide $\varepsilon$ añade
otro término, menor, del orden de $(p/R)\lambda T \times \varepsilon_0$, aquí 0.02 unidades.
\end{solution}

\begin{solution}{exo:b2:isotope-geochemistry:11}
$K = 13.3/5.0 = 2.66$. Con $s = 0.5$: $s' = 2.66\times0.5/(1 + 1.66\times0.5) =
1.33/1.83 = 0.73$. La diferencia $s' - s$ es máxima donde
$\mathrm{d}s'/\mathrm{d}s = K/(1+(K-1)s)^2 = 1$, en $s = (\sqrt K - 1)/(K-1) =
0.631/1.66 = 0.38$, donde $s' = 0.62$: una diferencia de 0.24.
\end{solution}

\begin{solution}{exo:b2:isotope-geochemistry:12}
Generación: $2.3$ a $2.6\times\qty{7.2e9}{km^3}/\qty{4.5e9}{yr} = 3.7$ a
\qty{4.2}{km^3} al año. Destrucción: $1.3$ a $1.6\times7.2/4.5 = 2.1$ a
\qty{2.6}{km^3} al año. Los ritmos actuales difieren: las adiciones (2.7) están por debajo de la
generación media del modelo, y la retirada (3.2) por encima de su destrucción media. En el modelo la corteza creció
pronto, cuando la generación era más rápida, y el presente es una época de equilibrio o de ligera
pérdida, como dice el balance de las zonas de subducción.
\end{solution}

\begin{solution}{pb:b2:isotope-geochemistry:1}
\textbf{1.} Manto: $(0.51307/0.512638 - 1)\times10^4 = +8.4$. Sedimento:
$(0.51141/0.512638 - 1)\times10^4 = -24.0$.

\textbf{2.} Kora $+2.8$, Vel $-4.3$, Asha $+4.7$.

\textbf{3.} Con $^{87}\mathrm{Sr}/^{86}\mathrm{Sr}$ en abscisas y $\varepsilon_{\mathrm{Nd}}$ en ordenadas:
manto (0.70282, $+8.4$), Kora (0.70410, $+2.8$), Vel (0.70618, $-4.3$), sedimento
(0.72000, $-24.0$), Asha (0.70290, $+4.7$). Solo Vel, con una
$\varepsilon_{\mathrm{Nd}}$ negativa, está en el lado enriquecido del CHUR.

\textbf{4.} Por debajo de la alineación: a la razón del estroncio del \omterm{def:b2:isotope-geochemistry:reservoirs}{manto empobrecido}, la alineación
tendría una $\varepsilon_{\mathrm{Nd}}$ cercana a $+8$; Asha tiene $+4.7$. Se parece al
grupo de Santa Elena, de tipo HIMU (Mangaia, Tubuai, Santa Elena).

\textbf{5.} La fusión cambia las razones padre--hijo pero no, de forma medible, las
razones isotópicas del estroncio y del neodimio (proposición sobre la fusión); las lavas son
jóvenes, de modo que sus razones no han crecido desde la erupción (pregunta 6).

\textbf{6.} $\Delta R = 0.12\,(e^{\lambda\cdot10^6} - 1) = 0.12\times\num{6.52e-6} =
\num{7.8e-7}$, es decir $\num{7.8e-7}/0.512638\times10^4 = 0.015$ unidades épsilon: muy por debajo de las
diferencias entre las islas, y por debajo de la precisión de una medida.

\textbf{7.} $\mathrm{Nd} = 40f + 0.75(1-f)$, y
$R = \bigl(40f\times0.51141 + 0.75(1-f)\times0.51307\bigr)/\mathrm{Nd}$.

\textbf{8.} $f = 0.5$ por ciento: Nd \qty{0.9463}{ppm}, $\varepsilon = +1.6$; 1 por ciento:
1.1425, $-2.9$; 2 por ciento: 1.535, $-8.4$; 5 por ciento: 2.7125, $-15.4$.

\textbf{9.} Sr 10.95, 11.90, 13.80 y \qty{19.50}{ppm}; $^{87}\mathrm{Sr}/^{86}\mathrm{Sr} =
0.70439$, 0.70571, 0.70780 y 0.71163.

\textbf{10.} $R\,\mathrm{Sr} = 10\times0.70282 + f(200\times0.72000 - 10\times0.70282)$ y
$\mathrm{Sr} = 10 + 190f$; eliminando $f$, $R\,\mathrm{Sr} = 7.0282 +
(\mathrm{Sr}-10)\times136.9718/190$, de modo que $R = 0.720904 - 0.18084/\mathrm{Sr}$: una recta
de pendiente $-0.1808$ (ppm) frente a $1/\mathrm{Sr}$. Comprobación en $f = 1$ por ciento:
$0.720904 - 0.18084/11.90 = 0.70571$.

\textbf{11.} $K = (10/0.75)/(200/40) = 13.33/5.0 = 2.67$. El manto es relativamente
más rico en estroncio, y el sedimento en neodimio, de modo que la razón del Nd se acerca más deprisa al
sedimento: un poco de sedimento desplaza mucho $\varepsilon_{\mathrm{Nd}}$ antes de que
$^{87}\mathrm{Sr}/^{86}\mathrm{Sr}$ se mueva apreciablemente.

\textbf{12.} 53 veces más rico en Nd ($40/0.75$), 20 veces en Sr. Con $f = 1$ por ciento el
sedimento aporta $0.40/1.1425 = 35$ por ciento del neodimio, lo que desplaza
$\varepsilon_{\mathrm{Nd}}$ en $0.35\times(8.4 + 24.0) = 11.3$ unidades, pero solo $2.0/11.9 =
17$ por ciento del estroncio.

\textbf{13.} $0.01\times40/1.1425 = 0.350$: un 35 por ciento.

\textbf{14.} A partir de $R\,(40f + 0.75(1-f)) = 40fR_S + 0.75(1-f)R_D$:
$f = 0.75(R - R_D)/\bigl(40(R_S - R) + 0.75(R - R_D)\bigr)$, y lo mismo con 200
y 10 para el estroncio.

\textbf{15.} $f = 0.75\times(-0.00029)/\bigl(40\times(-0.00137) + 0.75\times(-0.00029)\bigr)
= \num{-2.175e-4}/\num{-5.502e-2} = 0.0040$: un 0.40 por ciento.

\textbf{16.} $f = 10\times0.00128/\bigl(200\times0.01590 + 10\times0.00128\bigr) =
0.0128/3.1928 = 0.0040$: un 0.40 por ciento; los dos coinciden.

\textbf{17.} Neodimio: $0.75\times(-0.00065)/\bigl(40\times(-0.00101) +
0.75\times(-0.00065)\bigr) = \num{-4.875e-4}/\num{-4.089e-2} = 0.0119$; estroncio:
$10\times0.00336/\bigl(200\times0.01382 + 10\times0.00336\bigr) = 0.0336/2.7976 = 0.0120$.
Ambos dan un 1.2 por ciento.

\textbf{18.} $f = 0.75\times(-0.00019)/\bigl(40\times(-0.00147) + 0.75\times(-0.00019)\bigr)
= 0.0024$, que predice Sr \qty{10.46}{ppm} y $^{87}\mathrm{Sr}/^{86}\mathrm{Sr} =
0.70361$, frente a 0.70290 medido. Asha no está sobre la curva de mezcla manto--sedimento:
su fuente contiene otro componente.

\textbf{19.} Corteza oceánica reciclada antigua, el tipo HIMU: sus isótopos del Nd exigen una Sm/Nd
algo inferior a la del \omterm{def:b2:isotope-geochemistry:reservoirs}{manto empobrecido}, sus isótopos del Sr una Rb/Sr baja, y su
plomo (no medido aquí) una U/Pb alta, todo ello durante uno o dos mil millones de años.

\textbf{20.} $e^{\lambda_{147}\cdot2\times10^9} - 1 = 0.013133$. Sedimento: $0.51141 -
0.110\times0.013133 = 0.509965$; CHUR: $0.512638 - 0.1966\times0.013133 = 0.510056$;
$\varepsilon_{\mathrm{Nd}}(2\,\mathrm{Ga}) = (0.509965/0.510056 - 1)\times10^4 = -1.8$.

\textbf{21.} $f = 0.110/0.1966 - 1 = -0.440$; desde la subducción $\varepsilon$ ha cambiado en
$25.0\times(-0.440)\times2.0 = -22.0$, de modo que $\varepsilon(2\,\mathrm{Ga}) \approx -24.0 + 22.0
= -2.0$, a 0.2 unidades del valor exacto.

\textbf{22.} $e^{\lambda_{87}\cdot2\times10^9} - 1 = 0.028338$;
$0.72000 - 0.50\times0.028338 = 0.70583$.

\textbf{23.} En la subducción el sedimento tenía $\varepsilon_{\mathrm{Nd}} = -1.8$, cerca del
$-4.7$ de una corteza que entonces tenía solo 700 millones de años, y muy por encima de los
valores muy negativos de una corteza antigua: encaja con la erosión de una corteza joven, quizá
mezclada con algo de material derivado del manto. Su $^{87}\mathrm{Sr}/^{86}\mathrm{Sr}$ de
0.7058 en ese momento es también la de una roca cortical joven.

\textbf{24.} La fuente de Vel, la isla más enriquecida, contiene
\textbf{\qty[text-series-to-math=true]{1.2}{\percent}} de sedimento reciclado.
\end{solution}
